\section{Finite rational volume without measure theory}\label{app:volume}

This appendix supplies the elementary finite-volume
facts used in \cref{sec:polytope}, keeping that
reduction within arithmetic instead of appealing
to arbitrary Lebesgue measurable sets.

A rational polytope is a bounded solution set in
\(\Q^n\) of finitely many rational weak linear
inequalities. Explicit bounding inequalities can
be included. Faces, affine ranks, vertices, and
simplices are finite rational data.
Full-dimensional volume will be a rational number.

\subsection{Vertices can be enumerated constructively}

Suppose \(P=\{x:a_i\cdot x\leq b_i,\ i<h\}\)
has explicit box bounds. Enumerate every
\(n\)-element subset of rows. For each nonsingular
square system solve \(a_i\cdot x=b_i\) by rational
elimination and keep the solution if all
inequalities hold. Remove duplicates. Call this
finite list \(V\).

We prove by a bounded construction that every
rational \(x\in P\) is a convex combination of
vertices in \(V\). Inspect its tight rows.
If their rank is \(n\), an independent subset
already lists \(x\) in \(V\). Otherwise choose
a nonzero rational vector \(v\) in their common
nullspace.

Motion along \(v\) in each direction meets a
constraint, because of the explicit box bounds.
In the positive direction let
\[
 t_+=\min_{a_i\cdot v>0}
           \frac{b_i-a_i\cdot x}{a_i\cdot v}.
\]
Rows in this minimum are not tight, since \(v\)
annihilates every tight row, so \(t_+>0\).
The corresponding minimum for \(-v\) gives
\(t_->0\). Both endpoints
\(x_+=x+t_+v\) and \(x_-=x-t_-v\) lie in \(P\).
Each gains an independent tight row: the new row
does not annihilate \(v\), whereas the span of
the old tight rows does. Finally,
\[
 x=\frac{t_-}{t_++t_-}x_+
        +\frac{t_+}{t_++t_-}x_-.
\]
Repeat at the endpoints. Tight rank increases
on each branch, so after at most \(n\) stages
every leaf is in \(V\); there are at most
\(2^n\) leaves. Products of the branch weights
give a rational convex combination reconstructing
the original \(x\).

Conversely each vertex in \(V\) belongs to \(P\),
and its convex combinations satisfy the defining
inequalities. Thus \(P=\operatorname{conv}(V)\)
on rational points. Emptiness and affine dimension
are decided by finite vertex and rank calculations.
The argument applies to lower-dimensional
polytopes too: tight equations restrict the
allowable directions, and rank growth still
terminates at vertices of the bounded set.

\subsection{A compatible rational triangulation}

Enumerate nonempty faces by tight subsets of the
inequality list, identifying duplicates by their
vertex subsets. Let \(b_F\) be the average of the
vertices of each face \(F\).
An inequality not identically tight on \(F\) is
strict at at least one vertex, hence strict at
their average. Thus \(b_F\) is in the relative
interior of \(F\).

For every maximal flag
\[
 F_0\subset F_1\subset\cdots\subset F_d=F,
 \qquad \dim F_i=i,
\]
take the simplex with vertices
\(b_{F_0},\ldots,b_{F_d}\).
These simplices triangulate \(F\), compatibly
on all common faces.

In dimension zero this is immediate. Having
triangulated the facets compatibly, cone each
boundary simplex to \(b_F\).
For rational \(x\in F\), \(x\neq b_F\), the ray
\(b_F+t(x-b_F)\) reaches the boundary at the
least positive facet ratio, a finite minimum of
rational numbers. The resulting boundary point
and the construction in one lower dimension
give a simplex containing \(x\).
The first boundary point on the ray is unique;
together with disjoint relative interiors of
the boundary simplices, this proves disjoint
full interiors of the cones. Restrictions agree
on lower faces because the same face barycenters
were used throughout.

All recursion is over a finite face list with
strictly decreasing dimension. Rational affine
charts are obtained by selecting independent
vertex differences and solving rational systems.

\subsection{Simplex content and subdivision}

For a full-dimensional simplex define
\[
 c([v_0,\ldots,v_n])
   =\frac{|\det(v_1-v_0,\ldots,v_n-v_0)|}{n!}.
\]
Define \(\Vol(P)\) as the sum of these numbers in
the triangulation just constructed, and as zero
when \(\dim P<n\).

The determinant boundary identity is
\begin{equation}\label{eq:boundary}
 \det(v_1-v_0,\ldots,v_n-v_0)
 =\sum_{i=0}^n(-1)^i
      \det(v_0,\ldots,\widehat v_i,\ldots,v_n).
\end{equation}
Multilinearity proves it: terms with two copies
of \(v_0\) vanish; the remaining terms have
exactly the displayed signs.

Orient the small simplices of a triangulated
\(n\)-simplex consistently with the ambient
orientation and sum \cref{eq:boundary}.
Every interior \((n-1)\)-face occurs twice
with opposite induced orientations and cancels.
This incidence fact is finite: at a rational
relative-interior point of such a face, choose
a sufficiently small rational displacement to
either side. A minimum over the finitely many
other face inequalities supplies a permissible
size. Coverage gives an incident full simplex
on each side, and disjoint interiors give at
most one on each side.
On an exterior face, subdivision leaves its
determinant sum unchanged. To verify this last
claim, choose a rational affine chart of the
face. The determinant of its \(n\) vertex
position vectors is a fixed rational multiple
of the oriented \((n-1)\)-dimensional determinant
in the chart. If its hyperplane contains the
origin, the multiple is zero. Otherwise the
claim is the subdivision identity in dimension
\(n-1\). Dimension zero starts the induction.

Consequently the oriented determinant sum of
the small simplices equals that of the original
simplex. With positive orientation the
full-dimensional determinants are positive,
so absolute contents add.

Two finite rational triangulations have a finite
common refinement: intersect their simplices,
include the resulting cells and their faces,
and barycentrically triangulate this cell
complex. The subdivision identity applied to
each original simplex shows that both
triangulations have the same content.
It also proves additivity for a finite
face-to-face rational dissection.

\subsection{Monotonicity and covering}

Suppose rational polytopes \(Q_1,\ldots,Q_s\)
cover every rational point of a bounded
rational polytope \(P\). Cut a common bounding
box by all hyperplanes in their inequality
lists. The result is a finite rational cell
complex. The sign of each defining linear
form is constant in the relative interior of
a cell, and every full-dimensional cell has
a rational interior barycenter.

For a cell contributing to \(P\), its
barycenter is in some \(Q_j\) by the assumed
rational-point coverage. The constant sign
data put the entire cell in that \(Q_j\).
Nonnegativity and finite additivity give
\[
                    \Vol(P)\leq\sum_{j=1}^s\Vol(Q_j).
\]
Lower-dimensional cells have content zero.
Thus rational-point coverage suffices:
no sign or floor of an arbitrary real is
needed for this comparison.

\subsection{Boxes and affine transformations}

For invertible rational \(A\), determinant
multiplication gives
\[
 c([Av_0+b,\ldots,Av_n+b])
       =|\det A|\,c([v_0,\ldots,v_n]).
\]
A transformed triangulation remains rational.
Thus
\[
                  \Vol(AP+b)=|\det A|\,\Vol(P).
\]
Triangulate the unit cube into its \(n!\)
coordinate-order simplices. Each has absolute
determinant \(1\) and content \(1/n!\);
the cube's content is \(1\).
A diagonal affine map yields
\[
 \Vol\!\left(\prod_k[\ell_k,u_k]\right)
                         =\prod_k(u_k-\ell_k).
\]
Degenerate boxes have zero determinant and
content zero. In particular
\(\Vol(\lambda P)=\lambda^n\Vol(P)\)
for positive rational \(\lambda\).

\subsection{Primitive recursive bounds and induction}

Here are coarse explicit bounds. With \(h\)
inequalities there are at most
\(\binom hn\) candidate vertices,
\(2^h\) candidate faces, and
\((2^h)^{n+1}\) candidate flags.
A hyperplane refinement has at most \(3^h\)
sign patterns before infeasible cells are
discarded. A determinant can be computed by
its \(n!\) permutation expansion.
All rank, sign, incidence, minimum, and
equality tests are rational and decidable.

The splitting tree has depth at most \(n\).
The triangulation recursion visits only the
bounded face list. Numerator and denominator
sizes are bounded by iterating elementary
growth bounds for rational addition,
multiplication, and nonzero division through
the bounded number of operations.
This gives a primitive recursive bound on
the complete computation-state codes.

Correctness inductions can be expressed by
finite invariants. In vertex splitting, the
current list of nodes satisfies the inequalities;
its stored coefficients reconstruct the
original input point; weights are nonnegative
and sum to \(1\); unfinished nodes gain tight
rank at the next split. These are all bounded
Boolean equations. The rational point is a
numerical parameter, and its whole splitting
tree is finite. No induction over a statement
quantifying over every rational point is needed.

The triangulation and determinant arguments
likewise use a finite face-incidence list,
compatible boundary lists, and cancellation
of paired oriented faces.
Induction runs over the bounded lists of
faces and flags, storing the verified
determinant identities for earlier faces.
These are finite rational equalities.
Covering is checked on the finite list of
cell barycenters and sign patterns.
This gives the quantifier-free form of the
volume lemmas used in \cref{sec:polytope}.

The bound audit is part of the written proof.
A complete Lean realization of this appendix
remains a separate substantial implementation
task after the smaller certificate theorem.
